\documentclass[aspectratio=169,11pt]{beamer}
\usepackage[spanish]{babel}
\usepackage{fontspec}
\setsansfont{TeX Gyre Heros}
\usepackage{circuitikz}
\usepackage{booktabs}
\usepackage{siunitx}
\definecolor{azul}{HTML}{173B57}
\definecolor{morado}{HTML}{6A3D8F}
\setbeamercolor{structure}{fg=azul}
\setbeamercolor{frametitle}{fg=white,bg=azul}
\setbeamercolor{block title}{fg=white,bg=morado}
\setbeamercolor{block body}{fg=black,bg=morado!8}
\setbeamertemplate{navigation symbols}{}
\setbeamertemplate{footline}[frame number]
\title{Ley de Ohm: voltaje en una resistencia}
\subtitle{Laboratorio de Física Aplicada a las Ciencias Farmacéuticas}
\author{Diego García}
\date{}
\hypersetup{
  pdftitle={Ley de Ohm: voltaje en una resistencia},
  pdfauthor={Diego García}
}
\begin{document}

\begin{frame}
  \titlepage
\end{frame}

\begin{frame}{Prototipo de estructura}
\begin{alertblock}{No publicar}
Esta Experiencia 01 permite probar carpetas, documentos y carga no publicada
en Canvas. Los rangos y límites deben validarse con los equipos reales.
\end{alertblock}
\end{frame}

\begin{frame}{Pregunta experimental}
\centering
\Large ¿La caída de voltaje en una resistencia fija es proporcional a la
corriente que la atraviesa?

\vspace{1cm}
\[
V=RI
\]
\end{frame}

\begin{frame}{Qué vamos a medir}
\begin{columns}[T]
\column{0.48\textwidth}
\begin{itemize}
  \item \(R\) con el circuito apagado.
  \item \(I\) en serie.
  \item \(V\) en paralelo.
  \item Al menos seis puntos.
\end{itemize}
\column{0.48\textwidth}
\begin{circuitikz}[scale=0.8,transform shape,american currents]
  \draw (0,0) to[battery1] (0,3)
        to[ammeter,l={$A$}] (3,3)
        to[R,l={$R$}] (3,0) -- (0,0);
  \draw (3,3) -- (5,3)
        to[voltmeter,l={$V$}] (5,0) -- (3,0);
\end{circuitikz}
\end{columns}
\end{frame}

\begin{frame}{Instrumentos}
\begin{itemize}
  \item Fuente DC regulada con límite de corriente.
  \item Multímetro como voltímetro.
  \item Multímetro como amperímetro.
  \item Función ohmímetro con circuito desenergizado.
  \item Resistencia fija, placa y conductores.
\end{itemize}
\vfill
\textbf{Pendiente:} confirmar modelos, rangos, resolución, fusibles,
tolerancia y potencia.
\end{frame}

\begin{frame}{Tres reglas que protegen el equipo}
\begin{enumerate}
  \item El amperímetro nunca se conecta en paralelo con la fuente.
  \item El ohmímetro nunca se usa con el circuito energizado.
  \item Toda reconexión se realiza con la fuente apagada.
\end{enumerate}
\end{frame}

\begin{frame}{De los datos a la resistencia}
\begin{columns}
\column{0.48\textwidth}
\[
V=mI+b
\]
\begin{itemize}
  \item \(m\): resistencia estimada.
  \item \(b\): intercepto que debe discutirse.
  \item Comparar con ohmímetro y valor nominal.
\end{itemize}
\column{0.48\textwidth}
\begin{tikzpicture}[scale=0.85]
  \draw[->] (0,0)--(5,0) node[right] {\(I\)};
  \draw[->] (0,0)--(0,3.5) node[above] {\(V\)};
  \draw[azul,thick] (0.3,0.25)--(4.6,3.1);
  \foreach \x/\y in {0.7/0.55,1.4/1.0,2.1/1.48,2.8/1.85,3.5/2.45,4.2/2.8}
    \fill[morado] (\x,\y) circle (2pt);
\end{tikzpicture}
\end{columns}
\end{frame}

\begin{frame}{Taller: puente químico-farmacéutico}
\begin{block}{De una resistencia a una solución}
Una solución electrolítica también permite transporte de carga. Un
conductímetro relaciona respuesta eléctrica y conductancia, pero incorpora
geometría de celda, temperatura y calibración.
\end{block}
\vspace{0.4cm}
\textbf{Pregunta:} ¿por qué la resistencia metálica del montaje no basta para
describir una solución farmacéutica?
\end{frame}

\begin{frame}{A trabajar}
\begin{enumerate}
  \item Asignen cuatro roles.
  \item Midan \(R\) con la fuente apagada.
  \item Armen y soliciten revisión.
  \item Registren seis pares \(I,V\).
  \item Grafiquen, ajusten y concluyan.
\end{enumerate}
\end{frame}

\end{document}
